\section{Recurrent Neural Network}
\label{sec:rnn}

% TODO: fill this section once RNN experiments are complete (see 6_rnn.ipynb)

\subsection{Architecture}

% TODO: describe the RNN architecture.
%
% Key points to cover:
%   - Hidden state size, number of recurrent layers
%   - How temporal ordering of lag features is exploited (vs FNN flat vector)
%   - Activation: tanh (standard RNN cell)
%   - Dropout / regularization
%   - Output layer

\subsection{Training Setup}

% TODO: describe training procedure (same template as FNN section).

\subsection{Results}

% TODO: insert results table (MAE — full, calm, crisis).

% TODO: figure placeholder — predicted vs actual (RNN)

% TODO: figure placeholder — training / validation loss curves

\subsection{Limitations}

% TODO: describe remaining limitations of vanilla RNN:
%   1. Vanishing gradient — long-range dependencies are not captured.
%   2. Still struggles with abrupt regime shifts.
%
% These motivate LSTM in Section~\ref{sec:lstm}.
